% Generated draft: requires mathematical and source review.
\documentclass[11pt]{article}
\usepackage[margin=25mm]{geometry}
\usepackage{fontspec}
\setmainfont{Latin Modern Roman}
\usepackage{amsmath,amssymb,mathtools,unicode-math}
\setmathfont{Latin Modern Math}
\usepackage{enumitem,needspace,xurl,hyperref}
\setlength{\parindent}{0pt}
\setlength{\parskip}{5pt}
\setlength{\emergencystretch}{3em}
\newcommand{\N}{\mathbb{N}}
\newcommand{\Z}{\mathbb{Z}}
\newcommand{\Q}{\mathbb{Q}}
\newcommand{\R}{\mathbb{R}}
\newcommand{\C}{\mathbb{C}}
\newcommand{\F}{\mathbb{F}}
\newcommand{\A}{\mathbb{A}}
\newcommand{\PP}{\mathbb{P}}
\newcommand{\E}{\mathbb{E}}
\newcommand{\eps}{\varepsilon}
\newcommand{\dd}{\,\mathrm{d}}
\begin{document}
\section*{UnsolvedMath ID: 30005558}
\textbf{Source problem code:} catalogue entry
\subsection{Definitions and mathematical statement}
For integers \(d\geq 2\) and \(n\geq 1\), write
\[
\mathcal A_{n,d}=\operatorname{Adm}_{1}^{n+1}\bigl((2)^{2n}\bigr)_d
\]
for the source's notation for a moduli space of degree-\(d\) admissible covers of a genus-one curve with \(2n\) simple branch points and domain of genus \(n+1\). The supplied text explicitly explains this interpretation for \(d=2\); its higher-degree question uses the parallel notation. On the smooth connected locus, with no other ramification, the genus follows from Riemann--Hurwitz:
\[
2g(C)-2=d\bigl(2g(E)-2\bigr)+2n=2n.
\]
Connectedness, the treatment of boundary covers, and the precise moduli convention must be checked in the original source; they are not fully specified in the supplied extraction.

Under the usual interpretation, let \(\mathbb E\) be the rank-\(n+1\) Hodge bundle of the domain curve, whose fiber at a smooth domain \(C\) is \(H^0(C,\omega_C)\), and put \(\lambda_j=c_j(\mathbb E)\), with \(\lambda_0=1\). The interleaved source extraction appears to display the degree-two identity
\[
\sum_{n=1}^{\infty}\frac{u^{2n-1}}{(2n-1)!}
\int_{\mathcal A_{n,2}}\lambda_{n+1}\lambda_{n-1}
=\frac{i}{24}\,\frac{1-e^{iu}}{1+e^{iu}},
\qquad i^2=-1.
\]
This is a reconstruction of the supplied display, whose transcription and normalization require source review. The source attributes this example to joint work with H.-H. Tseng.

The natural reconstruction of Question 2 is: for every integer \(d>2\), determine a closed formula for
\[
F_d(u)=\sum_{n=1}^{\infty}\frac{u^{2n-1}}{(2n-1)!}
\int_{\mathcal A_{n,d}}\lambda_{n+1}\lambda_{n-1},
\]
and identify the rational function \(R_d\) satisfying
\[
F_d(u)=R_d(q),\qquad q=-e^{iu}.
\]
The source asserts that the parallel higher-degree expression is rational in \(q\). This assertion is supplied source information, not independently verified here. The displayed definition of \(F_d\) interprets ``parallel integral'' by retaining the degree-two Hodge classes and generating-series weights; the extracted question does not independently specify these choices for higher degree.

A completely precise formulation remains conditional on identifying the intended admissible-cover stack, its compactification and integration class, branch-point labeling and automorphism weights, whether the genus-one target varies, and the coefficient field and intended expansion or analytic meaning of the rational-function identity. These choices must not be inferred from the notation alone.
\subsection{Short English statement}
For each integer \(d>2\), the question asks for the explicit rational function of \(q=-e^{iu}\) represented by the higher-degree analogue of a supplied degree-two Hodge-integral series. The supplied record labels the problem open and asserts rationality, but the extracted text leaves the precise moduli and normalization conventions unresolved. This is a statement-only reconstruction requiring human mathematical review.
\subsection{Sources}
\begin{itemize}
\item Supplied local UnsolvedMath registry record, permanent ID \(30005558\), source problem code \texttt{catalogue entry}, title ``Hodge Integrals over Genus-One Admissible Covers,'' in the Oberwolfach Reports Open Problems collection. The record gives proposed year 2023 and no explicit proposer in its proposer field.
\item \emph{Recent Trends in Algebraic Geometry} (2023), Oberwolfach Reports, pp.~1485--1542, DOI \texttt{10.4171/owr/2023/27}. Supplied report reference: \url{https://doi.org/10.4171/owr/2023/27}.
\item Supplied literature-source attribution: Pandharipande, ``Maps to a moving elliptic curve,'' \url{https://doi.org/10.4171/owr/2023/27}. The supplied extract attributes the degree-two example to joint work with H.-H. Tseng.
\item Literature assessment dated exactly \texttt{2026-08-22}: the supplied curation labels the problem open, reports that the higher-degree integral is known to be rational in \(q\), and states that no authoritative general formula resolving the extracted question was located. This is supplied dated curation, not an independently verified literature assessment.
\item The registry labels its statement \texttt{corrected\_verified} and reports reconstruction and checking against the official report. These are supplied verification claims; no independent source visit or verification is asserted in this draft.
\end{itemize}
\end{document}
